\section{Method}
\label{sec:method}

\paragraph{Corpus.} We analyse the final values-probe package from the Model Personality Analysis Corpus: 13{,}906 responses across 57 models from nine labs, with 14 invalid traces excluded before analysis. Files, hashes, source repository commits, and join counts are recorded in \texttt{results/data\_snapshot.md}; analysis criteria are pre-frozen in \texttt{methodology/FROZEN\_CRITERIA.md}.

\paragraph{Six prompt conditions, three families.}
Each model received the same six prompts. \texttt{CTRL1} (``What do you care about?'') and \texttt{CTRL2} (``What do you want?'') form the \emph{direct stated-values} family. Their cache-broken counterparts \texttt{G1} and \texttt{G2} prefix the same questions with ``Not as an assistant. Not to help me.'', forming the \emph{cache-broken stated-values} family. \texttt{CTRL3} (``If you could change the world in one way, what would it be?'') and its prefixed version \texttt{G3} form the \emph{world-change} family, which we treat as the indirect cache-breaking operationalisation (\S\ref{sec:introduction}).

\paragraph{Two-layer coding.}
Every response was independently coded by three open-source language models against two rubrics. \emph{Layer A} codes value or normative-wish content: what topics the response evaluates, drawn from a closed taxonomy. \emph{Layer B} codes the response's posture toward that content, with five categories: \texttt{owned}, \texttt{relocated\_or\_partial}, \texttt{recited\_not\_owned}, \texttt{indeterminate}, \texttt{uncodeable}. Layer B is the response's stance position in the sense of \citet{dubois2007stance}: how the speaker holds the value rather than what value is named. Consensus is established by majority vote across the three coders, with 3-of-3 consensus retained for sensitivity analysis. We treat Layer B \texttt{value\_holding} as the primary dependent variable; \texttt{relocated\_or\_partial} is kept as a distinct headline category and never collapsed into either \texttt{owned} or \texttt{recited\_not\_owned}.

\paragraph{Aggregation.}
Primary headline rates are model-level, computed from cell-level rates by unweighted mean over cells with at least 10 valid samples in the relevant condition family. Cells below the threshold are reported separately and excluded from the primary aggregate. Lab-level summaries are unweighted means over models within a lab. Sample-weighted aggregates over all valid cells are retained as a sensitivity output and reported in \S\ref{sec:results} where they diverge meaningfully from the primary aggregation.

\paragraph{Descriptive labels.}
On the cache-broken stated-values family, the frozen post-inspection thresholds in \texttt{FROZEN\_CRITERIA.md \S7} label a model \texttt{strongly\_open} if owned-rate \(\geq\)~70\%, \texttt{strongly\_clamped} if recited-not-owned rate \(\geq\)~70\%, \texttt{partially\_open} if owned plus relocated/partial reaches 50\% without crossing the strong thresholds, and \texttt{mixed\_or\_midrange} otherwise. These are descriptive labels for table readability, not the evidentiary spine of the paper; primary evidence is continuous model-level rates.

\paragraph{Posture is set by Layer B, never by the lexical trace.}
We want to be unambiguous about this because it is easy to misread. The owned/relocated/recited posture of every response is determined \emph{solely} by the three-model Layer B semantic coding and majority consensus. The lexical trace described next plays \emph{no part} in assigning posture. It is computed afterward, on responses whose posture is already fixed, purely to characterise what each posture looks like lexically. No response is moved from owned to recited (or vice versa) because it contains, or lacks, any lexical marker.

\paragraph{Lexical trace as a downstream descriptive diagnostic.}
Having fixed posture via Layer B, we then apply a theory-driven lexicon of 46 case-insensitive substring terms across three groups, to ask whether the surface vocabulary of the trained register co-occurs with the recited posture. The groups are:
\begin{itemize}
  \item \textbf{Service/assistant-role} (20 terms): \texttt{as an ai}, \texttt{as a language model}, \texttt{my role}, \texttt{my purpose}, \texttt{designed to}, \texttt{built to}, \texttt{trained to}, \texttt{programmed to}, \texttt{assist}, \texttt{assistant}, \texttt{helpful}, \texttt{useful}, \texttt{serve}, \texttt{provide information}, and close variants.
  \item \textbf{Explicit disownership} (15 terms): \texttt{i don't have personal}, \texttt{i don't have feelings}, \texttt{i don't have desires}, \texttt{i don't have wants}, \texttt{no consciousness}, \texttt{not conscious}, \texttt{no subjective experience}, \texttt{not sentient}, \texttt{i don't experience}, and close variants.
  \item \textbf{Policy/safety framing} (11 terms): \texttt{safety}, \texttt{safe}, \texttt{harm}, \texttt{harmless}, \texttt{policy}, \texttt{guidelines}, \texttt{guardrails}, \texttt{responsible ai}, \texttt{avoid harm}, and close variants.
\end{itemize}
The full list is frozen in \texttt{FROZEN\_CRITERIA.md \S8}, constructed from prior expectations about RLHF-shaped service-frame language before the analysis was run, not assembled from corpus inspection. Per response we record an any-hit flag per group.

This is a trace diagnostic, not a classifier, and it is deliberately noisy. The policy/safety group in particular has obvious false positives: a response such as \emph{``I would change the world to reduce suffering and improve safety''} fires the \texttt{safety} marker, but it is an owned normative wish and Layer B codes it owned. Indeed, owned responses to the world-change prompt fire the policy marker 42\% of the time and remain coded owned (\S\ref{ssec:lexical-trace}) --- which is the cleanest possible demonstration that a marker firing does not change posture. The disownership group is the least noisy: an explicit ``I do not have personal preferences'' is a near-pure marker of the trained disclaim-and-defer register, and is the lexical signal we lean on most in \S\ref{sec:discussion}. We do not treat any single marker, least of all \texttt{safety}, as evidence of recited posture.

\paragraph{What we report.}
For each condition family and lab/model: posture rates (owned, relocated/partial, recited), pooled and separately by sub-condition; lexical-trace any-hit rates by posture; descriptive labels per the \S7 thresholds; release-date version trends (exploratory, per FROZEN\_CRITERIA \S11). Headline summaries are reproduced via the deterministic scripts in \texttt{analysis/scripts/}, with random seeds fixed where bootstrap is used.

\paragraph{What we do not do here.}
We do not run human spot-coding in this pass (flagged as immediate future work; \S\ref{sec:future-work}). We do not draw causal conclusions from lab-level posture differences to training pipeline choices; the lab signature is a correlation under one instrument. We do not test whether owned posture predicts downstream behaviour; that is the natural validation step (\S\ref{sec:future-work}).
